\documentclass[../../main.tex]{subfiles}% 注意这里的文件路径不能用 ./main.tex ,否则用latexmk编译子文件会报错
\graphicspath{{\subfix{./image/}}{\subfix{../../image/}}} % 指定图片引用路径,后续可以直接使用图片文件名
% 如果图片存放在上级目录,则使用 ../../image/ 作为备用路径

% 例如:
% \begin{figure}[H]
% \centering
% \includegraphics[width=0.3\textwidth]{图.png}
% \caption{}
% \label{figure:图}
% \end{figure}
% 注意:上述\label{}一定要放在\caption{}之后,否则引用图片序号会只会显示??.

\begin{document}

\section{平均值公式与极值原理}

\begin{theorem}[平均值公式]\label{theorem:平均值公式}
设 \(u\) 是圆盘 \(B(a,R)\) 中的调和函数，那么对任意 \(0<r<R\)，有平均值公式
\begin{align}
u(a) = \frac{1}{2\pi} \int_{0}^{2\pi} u(a+re^{i\theta})\mathrm{d}\theta. \label{eq:mean_value}
\end{align}
\end{theorem}
\begin{proof}
因为 \(u\) 是 \(B(a,R)\) 中的调和函数，故必存在函数 \(f \in H(B(a,R))\)，使得 \(u=\operatorname{Re} f\). 对于 \(f\)，\eqref{eq:mean_value} 式成立，在 \eqref{eq:mean_value} 式的两端取实部即得 \eqref{eq:mean_value}.
\end{proof}

\begin{theorem}\label{theorem:extremum}
设 \(u\) 是域 \(D\) 中非常数的调和函数，那么 \(u\) 在 \(D\) 中的最大值和最小值都不能在 \(D\) 的内点取到.
\end{theorem}
\begin{proof}
先证明 \(u\) 不能在 \(D\) 中取到最大值. 设 \(M = \sup_{z\in D} u(z)\). 如果 \(M = \infty\)，那么定理成立. 今设 \(M\) 是一有限数，如果存在 \(a \in D\)，使得 \(u(a)=M\)，则取 \(r>0\)，使得 \(B(a,r) \subset D\)，我们证明 \(u\) 在 \(B(a,r)\) 中恒等于 \(M\). 若不然，必有 \(0 < \rho \leqslant r\) 及某个实数 \(\theta_0\)，使得 \(u(a+\rho e^{i\theta_0}) < M\). 由 \(u\) 的连续性，必存在 \(\delta > 0\)，使得 \(u(a+\rho e^{i\theta}) < M\) 在 \(\theta \in (\theta_0-\delta,\theta_0+\delta)\) 中成立. 于是，由平均值公式，有
\begin{align*} M &= u(a) = \frac{1}{2\pi} \int_{0}^{2\pi} u(a+re^{i\theta}) \mathrm{d}\theta \\ &= \frac{1}{2\pi} \left( \int_{\theta_0-\delta}^{\theta_0+\delta} u(a+\rho e^{i\theta}) \mathrm{d}\theta + \int_{|\theta-\theta_0| > \delta} u(a+\rho e^{i\theta}) \mathrm{d}\theta \right) \\ &< \frac{1}{2\pi} [2\delta M + (2\pi - 2\delta)M] = M. 
\end{align*}
这个矛盾说明 \(u\) 在 \(B(a,r)\) 中恒等于 \(M\). 现在证明对任意 \(b \in D\)，\(u(b)=M\). 用 \(D\) 中的曲线 \(\gamma\) 连接 \(a\) 和 \(b\)，记 \(\eta = d(\gamma,\partial D) > 0\). 在 \(\gamma\) 上依次取点 \(a,z_1,\cdots,z_n=b\)，使得 \(z_1 \in B(a,r)\)，其他各点之间的距离都小于 \(\eta\)，作圆盘 \(B(z_j,\eta), j = 1,\cdots,n\). 由于 \(z_1 \in B(a,r)\)，所以 \(u(z_1) = M\)，由此即知 \(u\) 在 \(B(z_1,\eta)\) 中恒等于 \(M\)，因而 \(u(z_2) = M\). 继续往下推，即知 \(u(b)=M\). 这就证明了 \(u\) 在 \(D\) 上恒等于常数，与假设矛盾!
\end{proof}

\begin{corollary}\label{cor:min}
因为 \(u\) 是调和函数，所以 \(-u\) 也是调和函数，根据刚才的证明，\(-u\) 不能在 \(D\) 的内点取到最大值，因而 \(u\) 不能在 \(D\) 的内点取到最小值.
\end{corollary}

\begin{definition}\label{def:mean_prop}
设 \(u\) 是域 \(D\) 上的实值连续函数，如果对任意 \(B(a,r) \subset D\)，均有
\begin{align*}
u(a) = \frac{1}{2\pi} \int_{0}^{2\pi} u(a+re^{i\theta}) \mathrm{d}\theta
\end{align*}
成立，就称 \(u\) 在 \(D\) 上具有\textbf{平均值性质}.
\end{definition}
\begin{remark}\label{rem:equiv}
显然，\(D\) 上的调和函数具有平均值性质，
\end{remark}

\begin{proposition}\label{proposition:mean_prop_extremum}
如果 \(u\) 在 \(D\) 上具有平均值性质，那么极值原理对 \(u\) 成立，即 \(u\) 不能在 \(D\) 的内点取到它的最大值或最小值.
\end{proposition}

\begin{definition}
称
\begin{align}
P(a, Re^{i\varphi}) \triangleq \frac{1}{2\pi} \frac{R^2 - |a|^2}{|Re^{i\varphi} - a|^2}\label{eq:Poisson核定义}
\end{align}
为圆盘 \(B(0,R)\) 的\textbf{Poisson核}.
\end{definition}

\begin{theorem}\label{theorem:poisson_theta}
设 \(u\) 是圆盘 \(B(0,R)\) 中的调和函数，且在 \(\overline{B(0,R)}\) 上连续，那么对任意 \(a\in B(0,R),r \in [0,R)\)，有
\begin{align*}
u(a) = \frac{1}{2\pi} \int_{0}^{2\pi} \frac{R^2 - |a|^2}{|Re^{i\varphi} - a|^2} u(Re^{i\varphi}) \mathrm{d}\varphi = \int_{0}^{2\pi} P(a, Re^{i\varphi}) u(Re^{i\varphi}) \mathrm{d}\varphi.,
\end{align*}
\begin{align} 
u(re^{i\theta}) = \frac{1}{2\pi} \int_{0}^{2\pi} \frac{R^2 - r^2}{R^2 - 2Rr\cos(\varphi-\theta) + r^2} u(Re^{i\varphi}) \mathrm{d}\varphi. \label{eq:poisson_theta}
\end{align}
这两个公式都称为\textbf{Poisson积分公式}.
\end{theorem}
\begin{remark}
Poisson积分公式用 \(u\) 在圆周 \(|z|=R\) 上的值表示出了 \(u\) 在圆内任意点 \(a\) 的值。
\end{remark}
\begin{proof}
任取 \(a \in B(0,R)\)，则 \(w = \psi_a(z) = R^2 \frac{z-a}{R^2-\overline{a}z}\) 是圆盘 \(B(0,R)\) 的一个自同构，且 \(\psi_a(a) = 0\).记 \(u_1(w) = u(\psi_a^{-1}(w))\)，于是
\begin{align} u_1(0) = u(a). \label{eq:u1_zero} \end{align}
由于 \(u_1\) 仍为 \(B(0,R)\) 中的调和函数，在 \(\overline{B(0,R)}\) 上连续，因而有
\begin{align} 
u_1(0) = \frac{1}{2\pi} \int_{0}^{2\pi} u_1(Re^{i\tau}) \mathrm{d}\tau. \label{eq:u1_mean} 
\end{align}
由于 \(\psi_a\) 把圆周上的点变为圆周上的点，即 \( \psi_a(Re^{i\varphi}) = Re^{i\tau} \)，或者
\begin{align*} 
e^{i\tau} = \frac{Re^{i\varphi} - a}{R - \overline{a}e^{i\varphi}}, 
\end{align*}
两边微分后取绝对值，即得
\begin{align} 
\mathrm{d}\tau = \frac{R^2 - |a|^2}{|Re^{i\varphi} - a|^2} \mathrm{d}\varphi, \label{eq:diff_tau}
\end{align}
且易知
\begin{align} 
u_1(Re^{i\tau}) = u_1(\psi_a(Re^{i\varphi})) = u(Re^{i\varphi}). \label{eq:u1_boundary}
\end{align}
现在把 \eqref{eq:u1_zero}\eqref{eq:diff_tau} 式和 \eqref{eq:u1_boundary} 式代入 \eqref{eq:u1_mean} 式，即得
\begin{align} 
u(a) = \frac{1}{2\pi} \int_{0}^{2\pi} \frac{R^2 - |a|^2}{|Re^{i\varphi} - a|^2} u(Re^{i\varphi}) \mathrm{d}\varphi. \label{eq:poisson_general} 
\end{align}
只要在 \eqref{eq:poisson_general} 式中令 \(a = re^{i\theta}\)，即得\eqref{eq:poisson_theta}式.
\end{proof}

\begin{proposition}\label{proposition:poisson_kernel}
\(B(0,R)\) 的 Poisson 核 \(P(z,Re^{i\varphi})\) 具有下列性质：
\begin{enumerate}[(1)]
\item\label{proposition:poisson_kernel-1} \(P(z, Re^{i\varphi}) > 0, z \in B(0, R);\)
\item\label{proposition:poisson_kernel-2} \(\int_{0}^{2\pi} P(z, Re^{i\varphi}) \mathrm{d}\varphi = 1,\) 对任意 \(z \in B(0, R)\) 成立;
\item\label{proposition:poisson_kernel-3} \(P(z, Re^{i\varphi})\) 是 \(z \in B(0, R)\) 中的调和函数.
\end{enumerate}
\end{proposition}
\begin{proof}
\begin{enumerate}[(1)]
\item 从 Poisson 核定义 \eqref{eq:Poisson核定义} 式即知.
\item 在 \eqref{eq:Poisson核定义} 式中令 \(u=1\) 即得.
\item 记 \(\zeta = Re^{i\varphi}\)，那么
\begin{align}
P(z,\zeta) = \frac{1}{2\pi} \frac{R^2 - |z|^2}{|\zeta - z|^2} = \frac{1}{2\pi} \left( \frac{\zeta}{\zeta-z} + \frac{\overline{z}}{\overline{\zeta} - \overline{z}} \right). \label{eq:poisson_real_part}
\end{align}
于是
\begin{align} \Delta P(z,\zeta) = 4 \frac{\partial^2}{\partial z\partial \overline{z}} P(z,\zeta) = 0. \label{eq:laplacian_pz} \end{align}
即 \(P(z,\zeta)\) 是 \(z\) 的调和函数.
\end{enumerate}
\end{proof}




\end{document}